\documentclass{article}

% ============================================================
% Deep Learning 2026 Course Project Template
% ============================================================
\usepackage[final]{project_nips_style} % course style (based on NeurIPS layout)

% ============================================================
% Packages (import what you need)
% ============================================================
\usepackage[utf8]{inputenc}
\usepackage[T1]{fontenc}

\usepackage{hyperref}   % hyperlinks
\usepackage{url}        % simple URL typesetting

\usepackage{amsmath, amsfonts} % math
\usepackage{graphicx}   % figures
\usepackage{booktabs}   % nice tables
\usepackage{microtype}  % typography
\usepackage{xcolor}     % colors (use sparingly)

% Optional but often useful:
% \usepackage{siunitx}  % units and numbers
% \usepackage{cleveref} % better cross-references (Figure -> Fig.)
% \usepackage{listings} % code listings (if you want code blocks)

% ============================================================
% Title Information
% ============================================================
\title{Deep Learning 2026: Course Project Title}

\author{
  First Author \\
  University of Bern \\
  \texttt{first.last@unibe.ch}
  \And
  Second Author \\
  University of Bern \\
  \texttt{first.last@unibe.ch}
  % \And
  % Third Author \\
  % University of Bern \\
  % \texttt{first.last@unibe.ch}
}

\begin{document}
\maketitle

% ============================================================
% Abstract (keep short)
% ============================================================
\begin{abstract}
Briefly: (1) what problem you solve, (2) what method you built, (3) the main quantitative result.
Keep it tight and self-contained.
\end{abstract}

\vspace{0.5em}

% ============================================================
\section{Introduction and Background}
% ============================================================
\textbf{What to include:}
problem statement, why it matters, and a focused related-work summary.
Avoid a broad survey.

\paragraph{Problem statement.}
Define the task, inputs/outputs, and evaluation.
Example: Given input $x \in \mathbb{R}^{d}$, predict label $y \in \{1,\dots,K\}$.

\paragraph{Related work.}
Cite key papers: e.g., \citet{goodfellow2016deep} for fundamentals, \citet{he2016resnet} for residual nets.

\paragraph{Contributions (recommended).}
State 2--4 bullet points with what you actually did.
\begin{itemize}
  \item Implemented baseline model and strong training pipeline.
  \item Proposed method X and ablated key components.
  \item Evaluated on dataset Y with metric Z.
\end{itemize}

% ============================================================
\section{Method}
% ============================================================
\textbf{Goal:} make it possible for someone else to reproduce what you did.

\subsection{Model}
Describe architecture and important design choices.

\paragraph{Example equation: softmax classifier.}
For logits $z = f_\theta(x)$, the predicted probability is
\begin{equation}
p_\theta(y=k \mid x) = \frac{\exp(z_k)}{\sum_{j=1}^K \exp(z_j)}.
\end{equation}

\paragraph{Example loss: cross-entropy.}
\begin{equation}
\mathcal{L}(\theta) = -\frac{1}{N}\sum_{i=1}^{N}\log p_\theta\bigl(y_i \mid x_i\bigr).
\end{equation}

\subsection{Training setup}
Provide enough detail to replicate.
\begin{itemize}
  \item Dataset: name, split, preprocessing, augmentations
  \item Optimizer + hyperparameters (LR, weight decay, batch size, epochs)
  \item Hardware (GPU type) and training time (roughly)
\end{itemize}

\paragraph{Referencing figures/tables.}
Use labels and references: see Figure~\ref{fig:architecture} and Table~\ref{tab:main_results}.

\subsection{Implementation details}
Mention anything that changes results materially (seeds, schedules, normalization, gradient clipping, EMA, etc.).
If you used tricks, state them.

% ============================================================
\section{Results}
% ============================================================
\textbf{Goal:} numbers first, then interpretation.

\subsection{Main results}
Put your headline results in one table.

\begin{table}[t]
\centering
\caption{Main results on \textit{DatasetName}. Report mean $\pm$ std over 3 seeds if feasible.}
\label{tab:main_results}
\begin{tabular}{lcc}
\toprule
Method & Accuracy (\%) & F1 (\%) \\
\midrule
Baseline & 0.0 & 0.0 \\
Ours (full) & 0.0 & 0.0 \\
Ours (ablated) & 0.0 & 0.0 \\
\bottomrule
\end{tabular}
\end{table}

\subsection{Ablations / analysis}
Explain *why* you think results look like they do.
Avoid “it works well” without evidence.

\begin{itemize}
  \item What changed vs baseline?
  \item Where does it fail? (show examples or error categories)
  \item Sensitivity to hyperparameters?
\end{itemize}

% ============================================================
% Figures: how to include
% ============================================================
\begin{figure}[t]
  \centering
  % Replace with your file, e.g. figures/arch.pdf or .png
  \fbox{\rule[-.5cm]{0cm}{3.2cm} \rule[-.5cm]{5.2cm}{0cm}}
  \caption{Example figure. Use \texttt{\textbackslash includegraphics[width=\textbackslash linewidth]\{...\}} in real reports.}
  \label{fig:architecture}
\end{figure}

% Example real usage:
% \begin{figure}[t]
%   \centering
%   \includegraphics[width=0.9\linewidth]{figures/architecture.pdf}
%   \caption{Model architecture.}
%   \label{fig:architecture}
% \end{figure}

% ============================================================
\section{Conclusion}
% ============================================================
Summarize what you did and what you learned.
Be explicit about limitations and what you would do next.

\begin{itemize}
  \item Key outcome (with a number)
  \item Main limitation (honest)
  \item Next step (specific)
\end{itemize}

% ============================================================
% References (do not count toward page limit)
% ============================================================
{\small
\bibliographystyle{plainnat}
\bibliography{references}
}

% ============================================================
\appendix
% ============================================================
\newpage
\section{Appendix (Optional)}
Extra details, more plots, hyperparameter tables, qualitative examples, etc.
Not used in grading unless explicitly stated.

\end{document}